% !TeX root = main_fps.tex
\section{Background and Related Work}
Particle-based simulation encompasses a broad family of methodologies including molecular dynamics (MD),
discrete element methods (DEM), smoothed particle hydrodynamics (SPH), material point methods (MPM),
dissipative particle dynamics (DPD), and numerous other specialized formulations. These methods employ
different interaction models and mathematical formulations to represent particle behavior across a
wide range of physical applications.

Molecular dynamics (MD) derives particle interactions from interatomic potential functions that describe attraction, repulsion, 
and bonding behavior. Forces are obtained from the gradient of the potential-energy function and subsequently integrated to 
determine particle motion. The resulting dynamics are therefore closely coupled to the chosen potential model and have been 
widely used to study molecular systems, phase transitions, and atomistic phenomena.

Discrete Element Methods (DEM) model particle interactions through prescribed contact laws that relate particle overlap, 
relative motion, and material properties to interaction forces. Normal and tangential contact forces are commonly augmented 
with damping, friction, cohesion, or rolling-resistance models to reproduce desired material behavior. Consequently, particle 
dynamics emerge from the assumptions embedded within the selected contact formulation. DEM has been widely applied to granular 
materials, powders, and particulate systems where contact interactions play a dominant role.

Smoothed Particle Hydrodynamics (SPH) represents fluids through particles that carry continuum properties such as density, 
pressure, and velocity. Particle interactions are evaluated through kernel interpolation schemes that approximate continuum 
fields and their derivatives. Consequently, particle motion is governed by continuum relationships expressed through the SPH 
formulation rather than direct particle-contact interactions.

Despite their differences, existing particle methodologies commonly derive particle interactions from potential functions, 
contact laws, continuum approximations, or specialized response models. Particle methodologies also differ in the manner in
which interactions are evaluated and state updates are performed. These can significantly influence parallel
execution characteristics, synchronization requirements, and computational efficiency.

The Field Particle System adopts the source-owned interaction model introduced by gPCD. Particle
interactions are evaluated from the perspective of the source particle, and state updates are restricted
to quantities owned by the evaluating particle. This ownership model eliminates shared-state
modification during pair interactions, reduces synchronization requirements, and naturally supports
parallel execution. Consequently, force evaluation, acceleration accumulation, and state evolution
remain locally owned by each particle throughout the simulation.
